\chapter{Formation under the global free energy}
\label{ch:global-formation}

The global free energy of \cref{sec:dyn-likelihood} is the objective in which
a block of agents has one generative model, and it is the objective in which
formation is posed here. This chapter adds to that model a structural latent,
a nested pair of partitions with a prior on a declared finite catalog, block
latents at two levels, and single-agent records, and it measures the
resulting construction against the standard \cref{inh:recursive-formation}
with the objects that \cref{ch:inheritance} fixes before any flow is
analyzed. The general-law statements come first: the structure factor of the
recognition is a posterior of the structure in a model whose data are the
expected child kernels (\cref{gf:structure-factor}), and for child kernels of
exponential type it depends on the agents only through summed block
statistics, so the second-level structure inherits exactly from the first in
the static sense. Dynamical inheritance, type closure, and the closed forms
are proved for conjugate linear-Gaussian children and are features of that
family. The verdict is that the construction is a two-stage selection with
exact inheritance: (F1) to (F4) and (F6a) hold, (F5) fails because the
label-preserving null reproduces the partition sequence, and (F6b) fails
because the union forms from every counterfactual handoff.

The structure posterior is Bayesian model selection over a finite catalog
\citep{KassRaftery1995}, the conjugate Gaussian hierarchy is that of
\citet{LindleySmith1972}, and the structural recognition restricted to point
masses gives the classification version of the EM algorithm
\citep{CeleuxGovaert1992,DempsterLairdRubin1977}, read as coordinate descent
of a free energy \citep{NealHinton1998}. Bayesian hierarchical clustering
\citep{HellerGhahramani2005} computes a related structure posterior for a
greedy tree on fixed data. What is new here is that the data of the
structure posterior are the current laws of interacting agents, that the
population dynamics is descent of one free energy in which the structure is
a recognition factor, and that the resulting selection is measured against
null flows, protocol classes, and counterfactual handoffs.

\section{A two-level partition model}
\label{gf:sec-model}

\definitionheading{Nested structures and catalog}{gf:catalog}
Let $[N]=\{1,\dots,N\}$ be the agents. A nested structure is a pair
$\Pi=(\Pi^1,\Pi^2)$ in which $\Pi^1$ is a partition of $[N]$ into first-level
blocks $B$ and $\Pi^2$ is a partition of the set of blocks of $\Pi^1$ into
second-level blocks $C$. Write $B_\Pi(i)$ for the first-level block of agent
$i$ and $C_\Pi(B)$ for the second-level block of $B$. The finest structure
$\Pi_\perp$ has singleton blocks at both levels. A catalog is a finite set
$\mathsf C$ of nested structures containing $\Pi_\perp$, and the structure
prior is a law $\pi$ on $\mathsf C$ with $\pi(\Pi)>0$ for every
$\Pi\in\mathsf C$. Two structures induce the same agent partitions when their
first-level partitions agree and their second-level blocks contain the same
agents.

\constructionheading{Global model with a structural latent}{gf:model}
Let $\mathsf Y_i$, $\mathsf Z$, $\mathsf W$ be standard Borel spaces with a
sigma-finite reference $\lambda_i$ on $\mathsf Y_i$, let $r$ be a probability
law on $\mathsf W$, $k:\mathsf W\rightsquigarrow\mathsf Z$ a normalized
kernel, and $p_i(y\mid z)$ a jointly measurable probability density on
$\mathsf Y_i$ with respect to $\lambda_i$ for each $z\in\mathsf Z$. The
non-agent coordinate is $\nu=(\Pi,w,z)$ with $w=(w_C)_{C\in\Pi^2}$ and
$z=(z_B)_{B\in\Pi^1}$; it ranges over the standard Borel space
$\mathsf N=\bigsqcup_{\Pi\in\mathsf C}\mathsf W^{\Pi^2}\times\mathsf Z^{\Pi^1}$.
The prior is
\begin{equation}
 P_0(d\nu,dy)=\pi(\Pi)\prod_{C\in\Pi^2}r(dw_C)
 \prod_{B\in\Pi^1}k(w_{C_\Pi(B)},dz_B)
 \prod_{i=1}^Np_i(y_i\mid z_{B_\Pi(i)})\lambda_i(dy_i),
 \label{gf:prior}
\end{equation}
and agent $i$ has a single-agent record $o_i$ with a declared density
$\ell_i(o_i\mid y_i)$, positive and jointly measurable. The kernels $r$, $k$
and $p_i$ do not depend on the block, which is the homogeneity hypothesis
used in \cref{gf:first-event}. The objective is the global free energy
\cref{eq:dyn-fixed-vfe} of $\mathbb P=P_0\prod_i\ell_i$ at the record
$o=(o_i)_i$,
\begin{equation}
 \mathcal F(Q;o)=\KL(Q\Vert P_0)-\E_Q\sum_i\log\ell_i(o_i\mid y_i).
 \label{gf:free-energy}
\end{equation}
The conjugate linear-Gaussian instance takes
$\mathsf Y_i=\mathsf Z=\mathsf W=\R$, $r=\mathcal N(0,\sigma^2)$,
$k(w,\cdot)=\mathcal N(w,\tau^2)$, $p_i(\cdot\mid z)=\mathcal N(z,s_i^2)$, and
$\ell_i(o_i\mid y_i)=\mathcal N(o_i;y_i,\rho_i^2)$, with positive constants.
Given $\Pi$, the contents are then Gaussian with mean zero and covariance
\begin{equation}
 C_\Pi=S+\tau^2M^1_\Pi+\sigma^2M^2_\Pi,
 \qquad S=\operatorname{diag}(s_i^2),
 \label{gf:covariance}
\end{equation}
where $M^1_\Pi$ and $M^2_\Pi$ are the co-membership matrices of the first-level
blocks and of the agent sets of the second-level blocks; put
$R=\operatorname{diag}(\rho_i^2)$.

Agents interact in this model only through the shared non-agent coordinates.
No comparison record of the form \cref{eq:sample-comparison} is included,
for two reasons. Selection does not need it, since the shared block latents
already couple the agents of a block. And without factors that read two
agents, the agent blocks are conditionally independent given $\nu$ under
every product recognition, so their simultaneous update equals every
sequential one and the dynamics below has no order dependence among agents;
the nulls of \cref{inh:null-flow} then delete exactly the sharing of $\nu$.
The records are part of the configuration: they are coordinates held fixed
by the dynamics, so the admitted class below varies them with the initial
laws. This is the global picture's reading of evidence, which differs from
one population to another; the consequence of holding them fixed instead is
stated with (F4).

\definitionheading{Recognition families}{gf:recognition}
The structured family $\mathcal Q_{\mathrm{st}}$ consists of
$Q=Q_\nu\otimes\bigotimes_iQ_i$ with $Q_i=a_i\lambda_i$ and $Q_\nu$ an
arbitrary law on $\mathsf N$. The committed family
$\mathcal Q_{\mathrm{cm}}\subset\mathcal Q_{\mathrm{st}}$ restricts $Q_\nu$ to
$\delta_\Pi\otimes Q_{w,z\mid\Pi}$, a point mass on one structure times an
arbitrary law of the latents of that structure. Recognition is a product
across agents, as in \cref{prop:dyn-global-local}, and a product between the
agents and the non-agent block; inside the non-agent block it is not
factorized, so that the structure and the block latents keep their
dependence, which a factorization $Q(\Pi)\otimes Q(w,z)$ would discard. For
a discrete structure a point mass has finite entropy, so
$\mathcal F(\delta_\Pi\otimes Q_{w,z\mid\Pi}\otimes Q_y;o)\geq-\log\mathbb P(o)$
is a valid bound.

\propositionheading{The structure factor}{gf:structure-factor}
Fix agent laws $Q_i=a_i\lambda_i$ with $\int|a_i\log a_i|d\lambda_i<\infty$
and $\E_{Q_i}|\log\ell_i(o_i\mid\cdot)|<\infty$, and suppose that
$\int|\log p_i(y\mid z)|Q_i(dy)<\infty$ for every $z$. Put
\begin{equation}
 g_i(z)=\exp\int\log p_i(y\mid z)Q_i(dy),
 \qquad
 Z_\Pi=\int\prod_{C\in\Pi^2}r(dw_C)\prod_{B\in\Pi^1}\int k(w_{C_\Pi(B)},dz_B)
 \prod_{i\in B}g_i(z_B),
 \label{gf:pseudo-evidence}
\end{equation}
and suppose $0<Z_\Pi<\infty$ for every $\Pi\in\mathsf C$. Then:
\begin{enumerate}
\item[(i)] Over $\mathcal Q_{\mathrm{st}}$, the block minimizer of
$\mathcal F$ in $Q_\nu$ is the law proportional to
\begin{equation*}
 \pi(\Pi)\prod_Cr(dw_C)\prod_Bk(w_{C(B)},dz_B)\prod_ig_i(z_{B(i)}),
\end{equation*}
and its structure marginal is
\begin{equation}
 Q^*(\Pi)=\frac{\pi(\Pi)e^{-F_\Pi}}{\sum_{\Pi'}\pi(\Pi')e^{-F_{\Pi'}}},
 \qquad F_\Pi=-\log Z_\Pi.
 \label{gf:structure-posterior}
\end{equation}
\item[(ii)] Over $\mathcal Q_{\mathrm{cm}}$, the block minimizer selects
$\hat\Pi=\arg\min_\Pi[-\log\pi(\Pi)+F_\Pi]$, the mode of $Q^*$, with the
conditional latent law of (i) given $\hat\Pi$; it is unique when the mode is.
\item[(iii)] If the child densities are of exponential type,
$p_i(y\mid z)=h_i(y)\exp\langle\theta(z),T_i(y)\rangle$ with measurable
$\theta:\mathsf Z\to\R^p$ and $T_i:\mathsf Y_i\to\R^p$ and $\E_{Q_i}|T_i|<\infty$,
then $F_\Pi+\sum_i\E_{Q_i}\log h_i$ is a function of the block statistics
\begin{equation}
 t_B=\sum_{i\in B}\E_{Q_i}T_i(y_i),\qquad B\in\Pi^1,
 \label{gf:block-statistic}
\end{equation}
alone, so $Q^*$, the mode $\hat\Pi$, and for every first-level partition $A$
the conditional law of $\Pi^2$ and $w$ given $\Pi^1=A$ depend on the agent
laws only through $(t_B)_{B\in A}$.
\end{enumerate}
\begin{proof}
Under the stated integrability the free energy splits as in the proof of
\cref{prop:dyn-global-local}:
$\mathcal F=\KL(Q_\nu\Vert P_0^\nu)-\E_{Q_\nu}\sum_i\log g_i(z_{B(i)})
+\sum_i\int a_i\log a_id\lambda_i-\sum_i\E_{Q_i}\log\ell_i(o_i\mid y_i)$, where
$P_0^\nu$ is the prior marginal of $\nu$ and Fubini's theorem has been used
for the child factors. The first two terms equal
$\KL(Q_\nu\Vert Q_\nu^*)-\log\sum_\Pi\pi(\Pi)Z_\Pi$, with $Q_\nu^*$ the law in
(i), which proves (i) because relative entropy vanishes only at equality.
On $\mathcal Q_{\mathrm{cm}}$ the same computation with $Q_\nu=\delta_\Pi\otimes
Q_{w,z\mid\Pi}$ gives $-\log\pi(\Pi)-\log Z_\Pi+\KL(Q_{w,z\mid\Pi}\Vert
Q^*_{w,z\mid\Pi})$, which is (ii). Under (iii),
$\prod_{i\in B}g_i(z)=\exp\big(\langle\theta(z),t_B\rangle+\sum_{i\in B}
\E_{Q_i}\log h_i\big)$, and every agent lies in exactly one first-level block
of every structure, so the factor $\exp\sum_i\E_{Q_i}\log h_i$ is common to
all $Z_\Pi$ and the rest of \cref{gf:pseudo-evidence} reads only $t_B$. The
conditional law of $(\Pi^2,w)$ given $\Pi^1=A$ is proportional to
$\pi(A,\Pi^2)\prod_Cr(dw_C)\prod_{B\in A}L_B(w_{C(B)})$ with
$L_B(w)=\int k(w,dz)e^{\langle\theta(z),t_B\rangle}$.
\end{proof}

Part (iii) is the general-law form of the statement that the second-level
posterior inherits from the first level through a block retention: given
$\Pi^1=A$ the second level sees a one-level model whose units are the blocks
of $A$ and whose child likelihoods are $L_B$. It needs no conjugacy of $k$
or $r$, which enters only the closed form of $Z_\Pi$. The structure factor is
a posterior of the structure, but in the model whose data are the expected
child kernels $g_i$, not in the model of the records: the two coincide only
in the limit described after \cref{gf:gaussian-structure}.

\corollaryheading{The Gaussian structure factor}{gf:gaussian-structure}
In the conjugate linear-Gaussian instance, let $Q_i$ have mean $m_i$ and
variance $S_i$. Then
\begin{equation}
 -\log\pi(\Pi)+F_\Pi
 =-\log\pi(\Pi)-\log\mathcal N(m;0,C_\Pi)+\sum_i\frac{S_i}{2s_i^2}
 +\mathrm{const},
 \label{gf:gaussian-score}
\end{equation}
with a constant independent of $\Pi$ and of $Q$, so $Q^*$ is the posterior of
the structure when the vector of current means $m$ is observed with child
noise $S$, and the selector
$\hat\Pi(m)=\arg\max_\Pi\pi(\Pi)\mathcal N(m;0,C_\Pi)$ is Bayesian model
selection on the means. The block statistic is
$t_B=(\lambda_B\hat m_B,\lambda_B)$ with
$\lambda_B=\sum_{i\in B}s_i^{-2}$ and
$\hat m_B=\lambda_B^{-1}\sum_{i\in B}s_i^{-2}m_i$.
\begin{proof}
Here $\log p_i(y\mid z)=-(y-z)^2/(2s_i^2)+\mathrm{const}$ and
$\E_{Q_i}(y-z)^2=(m_i-z)^2+S_i$, so $g_i(z)$ is proportional to
$\mathcal N(m_i;z,s_i^2)e^{-S_i/(2s_i^2)}$, and $Z_\Pi$ is the marginal
density of $m$ in the hierarchy \cref{gf:prior}, which is
$\mathcal N(m;0,C_\Pi)$. In the form of (iii), $\theta(z)=(z,-z^2/2)$ and
$T_i(y)=(y/s_i^2,1/s_i^2)$.
\end{proof}

The variance term in \cref{gf:gaussian-score} is the same for every
structure because the child noise of an agent does not depend on its block.
The selector therefore reads only the means, and a non-Gaussian initial law
with finite second moments gives the same selection as the Gaussian law with
its mean and variance.

\section{Dynamics, chart, and persistence}
\label{gf:sec-dynamics}

\definitionheading{The declared dynamics}{gf:sweep}
The dynamics is coordinate descent of \cref{gf:free-energy} on
$\mathcal Q_{\mathrm{cm}}$ in sweeps. Sweep $n\mapsto n+1$ replaces $Q_\nu$ by
its block minimizer given $Q_y^{(n)}$ (\cref{gf:structure-factor}(ii)), and
then every $Q_i$ by its block minimizer given the new $Q_\nu$,
\begin{equation}
 Q_i^{(n+1)}(dy)\propto\ell_i(o_i\mid y)\exp\Big(\int\log p_i(y\mid z)
 Q_\nu^{(n+1)}(dz_{B(i)})\Big)\lambda_i(dy).
 \label{gf:agent-update}
\end{equation}
The selector is the measurable map $x\mapsto\hat\Pi(x)$ of
\cref{gf:structure-factor}(ii) evaluated at the agent laws and the records of
a configuration $x$, on the set where the mode is unique; along the declared
dynamics the structural recognition of sweep $n+1$ equals $\hat\Pi(x_n)$.
The membership history is $H(x)=(\hat\Pi(x_n))_{n\geq0}$ with the sweep index
as time, and $P(x)$ is its sequence of distinct consecutive values.

This is the dynamics of \cref{prop:dyn-global-local}, each block being moved
to the minimizer of its local free energy, with the non-agent coordinate as
one more block. It is chosen over a continuous Fisher flow because every
step is an exact block minimization, so $\mathcal F$ is nonincreasing along
it, and because in the Gaussian instance every step is in closed form. The
committed family is chosen over $\mathcal Q_{\mathrm{st}}$ because the block
retention below is exact on it; on $\mathcal Q_{\mathrm{st}}$ the agent update
mixes the structures with weights $Q^*(\Pi)$ that depend on within-block
statistics through the structures with another first level, so that
retention is not exact there (\cref{gf:soft-remark}).

\propositionheading{The Gaussian sweep}{gf:gaussian-sweep}
In the conjugate linear-Gaussian instance, let $Q_i^{(0)}$ have finite second
moments. For every $n\geq1$, $Q_i^{(n)}=\mathcal N(m_i^{(n)},\eta_i)$ with
$\eta_i=(s_i^{-2}+\rho_i^{-2})^{-1}$, and with
$a_i=s_i^2/(s_i^2+\rho_i^2)$, $b_i=1-a_i$, $\mathsf a=\operatorname{diag}(a_i)$,
$\mathsf b=\operatorname{diag}(b_i)$,
\begin{equation}
 \Pi^{(n+1)}=\hat\Pi(m^{(n)}),\qquad
 m^{(n+1)}=\mathsf ao+\mathsf bK_{\Pi^{(n+1)}}m^{(n)},\qquad
 K_\Pi=I-SC_\Pi^{-1}.
 \label{gf:mean-map}
\end{equation}
The vector $K_\Pi m=m+S\nabla\log\mathcal N(m;0,C_\Pi)$ is the posterior mean of
the block latents $(z_{B(i)})_i$ given the means as data. With
$\bar F(m)=\sum_i(o_i-m_i)^2/(2\rho_i^2)-\log\mathcal N(m;0,C_\Pi)$, the mean
update is $m-\operatorname{diag}(\eta_i)\nabla\bar F(m)$, a unit step of the
natural gradient for the Fisher metric of the agent laws at their stationary
variances. For fixed $\Pi$ the map $m\mapsto\mathsf ao+\mathsf bK_\Pi m$ is a
contraction with factor at most $\max_ib_i$ in the norm
$\Vert x\Vert_S=(\sum_ix_i^2/s_i^2)^{1/2}$, and its fixed point is
\begin{equation}
 m^*_\Pi=C_\Pi(C_\Pi+R)^{-1}o=\E[y\mid o,\Pi],
 \qquad K_\Pi m^*_\Pi=(C_\Pi-S)(C_\Pi+R)^{-1}o,
 \label{gf:fixed-point}
\end{equation}
the exact posterior means of the contents and of the block latents given the
records and the structure.
\begin{proof}
Given $\Pi$ the latent law of \cref{gf:structure-factor}(i) is the Gaussian
posterior of $u=(z_{B(i)})_i$ given the data $m$ with noise $S$, whose mean is
$(C_\Pi-S)C_\Pi^{-1}m=K_\Pi m$; since $\nabla\log\mathcal N(m;0,C_\Pi)=
-C_\Pi^{-1}m$, this is the stated identity. In \cref{gf:agent-update} the
exponent is $-(y-\bar z_i)^2/(2s_i^2)$ up to a constant, with $\bar z_i=(K_\Pi
m)_i$, and the product with the record density is
$\mathcal N(\eta_i(\rho_i^{-2}o_i+s_i^{-2}\bar z_i),\eta_i)$, which gives
\cref{gf:mean-map}; the gradient form follows from
$\nabla\bar F=-R^{-1}(o-m)+C_\Pi^{-1}m$ and $\eta_i\rho_i^{-2}=a_i$,
$\eta_is_i^{-2}=b_i$. The matrix $S^{-1/2}K_\Pi S^{1/2}=I-S^{1/2}C_\Pi^{-1}S^{1/2}$
is symmetric with spectrum in $[0,1)$ because $C_\Pi\geq S>0$, and
$\mathsf b$ commutes with $S$, which gives the contraction. At the fixed
point, $(I-\mathsf bK_\Pi)m=\mathsf ao$ and $I-\mathsf bK_\Pi=\mathsf a
+\mathsf bSC_\Pi^{-1}=\mathsf a(I+RC_\Pi^{-1})$ because $\mathsf a^{-1}\mathsf
bS=R$, so $m^*_\Pi=(I+RC_\Pi^{-1})^{-1}o=C_\Pi(C_\Pi+R)^{-1}o$, and the second
identity follows from $K_\Pi C_\Pi=C_\Pi-S$.
\end{proof}

Equality of the mean-field means with the exact posterior means at fixed
structure, data-independent variances, and invariance of the Gaussian family
are features of the Gaussian family, as in \cref{prop:dyn-gaussian-closure};
for other child or record kernels \cref{gf:agent-update} leaves no
finite-dimensional family invariant in general. The stationary relation
between the selector and the posterior of the records is exact but not an
equality: a structure $\Pi_f$ at which the dynamics comes to rest satisfies
$\Pi_f=\hat\Pi(\E[y\mid o,\Pi_f])$, a self-consistency condition, while the
maximum a posteriori structure is $\arg\max_\Pi\pi(\Pi)\mathcal N(o;0,C_\Pi+R)$.
As $R\to0$ both tend to $\arg\max_\Pi\pi(\Pi)\mathcal N(o;0,C_\Pi)$, since
$m^*_\Pi\to o$. For $R>0$ the pseudo-data $m^*_\Pi$ have within-block
contrasts $a_i$ times those of the records but are read with noise $S$ rather
than $S+R$, so the selector can commit to a structure that the posterior of
the records does not prefer.

\lemmaheading{Persistence}{gf:persistence}
In the Gaussian instance, suppose $\Pi_f$ is the strict maximizer of
$\pi(\Pi)\mathcal N(m;0,C_\Pi)$ on the ball
$\{\Vert m-m^*_{\Pi_f}\Vert_S\leq r\}$ for some $r>0$. If $\Pi^{(n+1)}=\Pi_f$ and
$m^{(n)}$ lies in the ball, then $\Pi^{(k)}=\Pi_f$ for all $k>n$ and
$m^{(k)}\to m^*_{\Pi_f}$ geometrically.
\begin{proof}
The map with $\Pi_f$ fixed maps the ball into itself by the contraction of
\cref{gf:gaussian-sweep}, and on the ball the selector returns $\Pi_f$.
\end{proof}

The persistence criterion of (F2) is therefore an invariant block relation,
certified by a ball of this lemma.

\section{Inheritance, top-down action, and type closure}
\label{gf:sec-inheritance}

\theoremheading{Exact inheritance through the block retention}{gf:inheritance}
In the Gaussian instance, let $A$ be a first-level partition and let
$\Omega_A$ be a set of configurations, invariant under the sweep, at every
point of which $\hat\Pi$ has first level $A$ and the agent variances equal
$\eta_i$. For $B\in A$ put
\begin{equation}
 R_B(x)=\big(\lambda_B,\hat m_B,\hat o_B\big),\qquad
 \hat o_B=\frac{\sum_{i\in B}\alpha_io_i}{\sum_{i\in B}\alpha_i},
 \quad\alpha_i=\frac{\eta_i}{s_i^2\rho_i^2},
 \label{gf:retention}
\end{equation}
and let $R_A=(R_B)_{B\in A}$. Then on $\Omega_A$:
\begin{enumerate}
\item[(L)] the parent law $\mathcal N(\hat m_B,\eta_B)$, with
$\eta_B=\lambda_B^{-1}b_B$ and $b_B=\lambda_B^{-1}\sum_{i\in B}s_i^{-4}\eta_i$,
is the law of $\hat y_B=\lambda_B^{-1}\sum_{i\in B}s_i^{-2}y_i$ under $Q$, and
the second-level structure, the second-level latent law, and the block
latent laws are functions of $R_A$;
\item[(D)] $R_A\circ T=T^c\circ R_A$ for the sweep $T$, with the parent sweep
\begin{equation}
 \hat\Pi^2=\arg\max_{\Pi^2}\pi(A,\Pi^2)\mathcal N(\hat m;0,C^c_{\Pi^2}),\qquad
 \hat m_B'=a_B\hat o_B+b_B\big(K^c_{\hat\Pi^2}\hat m\big)_B,
 \label{gf:parent-sweep}
\end{equation}
where $C^c_{\Pi^2}=\operatorname{diag}(\lambda_B^{-1})+\tau^2I+\sigma^2M_{\Pi^2}$,
$K^c=I-\operatorname{diag}(\lambda_B^{-1})(C^c)^{-1}$ and $a_B=1-b_B$; it is
globally defined and well posed;
\item[(B)] for the protocol that averages one fresh sample from each member
with weights $s_i^{-2}/\lambda_B$, the response law is $\mathcal N(\hat
m_B,\eta_B)$, a function of $R_A$;
\item[(C)] the channels read by the parent update, the block likelihood
$z\mapsto\prod_{i\in B}\mathcal N(m_i;z,s_i^2)\propto\mathcal N(\hat
m_B;z,\lambda_B^{-1})$ and the parent record term, factor through $R_A$.
\end{enumerate}
\begin{proof}
For $\Pi=(A,\Pi^2)$ decompose $m$ into the block means $\hat m$ and the
within-block contrasts. The matrices $M^1_\Pi$ and $M^2_\Pi$ are constant on
the blocks of $A$, so the hierarchy reads $m$ only through
$\prod_{i\in B}\mathcal N(m_i;z_B,s_i^2)$, which is
$\mathcal N(\hat m_B;z_B,\lambda_B^{-1})$ times a factor that does not depend
on $z_B$ or on $\Pi^2$. Hence $\log\mathcal N(m;0,C_\Pi)=W_A(m)+\log\mathcal
N(\hat m;0,C^c_{\Pi^2})$ with $W_A$ independent of $\Pi^2$, and the latent
posteriors given $\Pi$ depend on $m$ through $\hat m$; this proves the second
half of (L) and the first part of \cref{gf:parent-sweep}. All agents of $B$
receive the same $\bar z_B=(K_\Pi m)_i$, so by \cref{gf:mean-map}
\[
 \lambda_B\hat m_B'=\sum_{i\in B}s_i^{-2}\eta_i\big(\rho_i^{-2}o_i
 +s_i^{-2}\bar z_B\big)
 =\Big(\sum_{i\in B}\alpha_i\Big)\hat o_B+\lambda_Bb_B\bar z_B,
\]
and $\alpha_i+s_i^{-4}\eta_i=s_i^{-2}$ gives $\sum_{i\in B}\alpha_i=\lambda_Ba_B$.
The block latent mean $\bar z_B$ is the parent's $(K^c\hat m)_B$, which
proves \cref{gf:parent-sweep}. The variance of $\hat y_B$ under $Q$ is
$\lambda_B^{-2}\sum_{i\in B}s_i^{-4}\eta_i=\eta_B$, and the protocol in (B)
has that law. The parent sweep is an affine contraction composed with a
selection over a finite set, hence defined for every parent state, and (C)
is the factorization just used.
\end{proof}

The theorem is the dynamical counterpart of \cref{gf:structure-factor}(iii)
and holds only on a set where the first level stays $A$: the full selector
compares structures with other first levels, whose scores read the
within-block contrasts, and the invariance of $\Omega_A$ is a property of the
trajectories that must be checked. The closure is a feature of conjugate
linear-Gaussian children with linear-Gaussian records: the block sum of
updated statistics is a function of $\hat o_B$ because the update is affine
in the record.

\propositionheading{What fails beyond the affine update}{gf:beyond}
Under \cref{gf:structure-factor}(iii), with $\Pi$ fixed, the update
\cref{gf:agent-update} gives $\E_{Q_i^{(n+1)}}T_i=G_i(o_i,\bar\theta_B)$, where
$\bar\theta_B=\E_{Q_\nu}\theta(z_B)$ and $G_i(o,\vartheta)$ is the gradient
in $\vartheta$ of $\log\int\ell_i(o\mid y)h_i(y)e^{\langle\vartheta,T_i(y)\rangle}
\lambda_i(dy)$, wherever this logarithm is finite and differentiable. If there are measurable record statistics
$\varsigma_i$ and maps $\Gamma_B$ with $\sum_{i\in B}G_i(o_i,\vartheta)=
\Gamma_B(\vartheta,\sum_{i\in B}\varsigma_i(o_i))$, then the retention
$(t_B,\sum_{i\in B}\varsigma_i(o_i))$ satisfies (D) wherever the first level
stays $A$. Conversely, for two agents with the same kernels, records in an
interval $J\subset\R$, $\varsigma$ the identity, and $G(\cdot,\vartheta)$
continuous, the condition forces $G(\cdot,\vartheta)$ to be affine on $J$.
\begin{proof}
The first statement is the computation in the proof of
\cref{gf:inheritance} with $t_B$ in place of $\lambda_B\hat m_B$, using
\cref{gf:structure-factor}(iii) for the structure and latent updates. For the
converse, $G(x)+G(y)=\Gamma(x+y)$ for $x,y\in J$ implies
$G(x)+G(y)=2G(\tfrac{x+y}2)$, since $(\tfrac{x+y}2,\tfrac{x+y}2)$ has the same
sum. A midpoint-affine function agrees with the affine interpolant of any
two of its values at all dyadic points between them, and continuity extends
the agreement to the interval.
\end{proof}

For a non-Gaussian record kernel the static inheritance of
\cref{gf:structure-factor}(iii) survives, while dynamical inheritance with a
finite block retention requires an update affine in a record statistic;
otherwise the retention must carry the individual records of the block,
which is the injective retention excluded by (F4). Whether an augmented
finite parent state exists for such families is open.

\remarkheading{The structured family}{gf:soft-remark}
On $\mathcal Q_{\mathrm{st}}$ the stationary $Q_\nu$ is the law of
\cref{gf:structure-factor}(i), and the Gaussian agent update uses
$\bar z=\sum_\Pi Q^*(\Pi)K_\Pi m=m+S\nabla\log\sum_\Pi\pi(\Pi)\mathcal
N(m;0,C_\Pi)$. A within-block contrast $\delta$ for $A$ leaves each
$\mathbf 1_B^{\mathsf T}K_\Pi(m+\delta)$ with $\Pi^1=A$ unchanged but changes
the weights $Q^*(\Pi)$ of the structures with another first level, which
have positive prior mass; the block sums of the next means then depend on
$\delta$, and the retention \cref{gf:retention} is not exact. It becomes
exact in the limit in which $Q^*$ concentrates on $\Pi^1=A$. The dynamics and
selection of the structured family are not analyzed here.

\propositionheading{Top-down action and the cavity prior}{gf:top-down}
In the model \cref{gf:model} with $\Pi$ fixed, the parent acts on child $i$
in \cref{gf:agent-update} through the factor
$\exp\int\log p_i(y\mid z)Q_\nu(dz_{B(i)})$, the geometric mean of the child
kernel under the block-latent law. The exact posterior marginal satisfies
\begin{equation}
 \mathbb P(dy_i\mid o,\Pi)\propto\ell_i(o_i\mid y_i)c_i(y_i)\lambda_i(dy_i),
 \qquad
 c_i(y_i)=\int p_i(y_i\mid z)\mathbb P(dz_{B(i)}\mid o_{-i},\Pi),
 \label{gf:cavity}
\end{equation}
with the cavity prior $c_i$, which uses the records of the other agents and
not that of $i$. In the Gaussian instance at the fixed point
\cref{gf:fixed-point}, the top-down factor is
$\mathcal N(y;\E[z_{B(i)}\mid o,\Pi],s_i^2)$, which uses $o_i$ and differs from
the cavity prior $\mathcal N(\bar c_i,s_i^2+v_i)$; the two give the same
posterior mean of $y_i$, and the mean-field variance $\eta_i$ is at most the
exact posterior variance $[C_\Pi-C_\Pi(C_\Pi+R)^{-1}C_\Pi]_{ii}$.
\begin{proof}
Given $\Pi$, $o_i$ depends on the other variables only through $y_i$, and
$y_i$ is conditionally independent of $o_{-i}$ given $z_{B(i)}$, so
$\mathbb P(dy_i\mid o,\Pi)\propto\ell_i(o_i\mid y_i)\mathbb P(dy_i\mid o_{-i},\Pi)$
and the second factor is $c_i\lambda_i$. The Gaussian means agree by
\cref{gf:fixed-point}. The variance inequality is the classical fact that a
mean-field factor of a Gaussian has the inverse diagonal precision as its
variance, here $(s_i^{-2}+\rho_i^{-2})^{-1}$, which is at most the diagonal of
the covariance \citep{Bishop2006}.
\end{proof}

The parent acts on its children through a term of the fine field that is
absent before the handoff: with $\Pi_\perp$ selected agent $i$ is drawn
toward its own latent, and after the handoff toward the shared block latent.
This top-down action is the mean-field one, not the cavity prior of exact
inference.

\propositionheading{Type closure}{gf:type-closure}
Let $\mathfrak G$ be the class of conjugate linear-Gaussian nested-structure
models of \cref{gf:model} with agent-dependent $s_i^2,\rho_i^2$, any catalog
and positive prior, single-agent records, and the sweep of \cref{gf:sweep}.
Under the hypotheses of \cref{gf:inheritance}, the parent sweep
\cref{gf:parent-sweep} is the sweep of the model of $\mathfrak G$ whose agents
are the blocks of $A$, with child noise $\lambda_B^{-1}$, record
$\hat o_B$ with noise $\rho_B^2=\lambda_B^{-1}b_B/a_B$, first-level catalog
the singletons, second-level catalog the partitions $\Pi^2$ with
$(A,\Pi^2)\in\mathsf C$, and prior proportional to $\pi(A,\cdot)$; its agent
laws $\mathcal N(\hat m_B,\eta_B)$ are its block minimizers. The parent is
type-closed for $\mathfrak G$ in the sense of \cref{inh:type-closed}, with the
coordinate-descent dynamics in place of a metric flow.
\begin{proof}
With these constants $a_B=\lambda_B^{-1}/(\lambda_B^{-1}+\rho_B^2)$, and the
parent model's update \cref{gf:mean-map} is \cref{gf:parent-sweep}; its
stationary variance is $(\lambda_B+\rho_B^{-2})^{-1}=\lambda_B^{-1}b_B=\eta_B$.
The structure score is \cref{gf:gaussian-score} for the parent model, which
is the second term of the decomposition in the proof of
\cref{gf:inheritance}.
\end{proof}

The parent's first level is forced to be the singleton catalog, so type
closure lowers the depth of the hierarchy by one; it does not make the
parent able to regroup its own children.

\section{Controls under the global free energy}
\label{gf:sec-controls}

The nulls and the frozen-environment control of the global free energy are
those of \cref{inh:null-flow}. In this model the label-preserving null gives
each agent a private copy of $(\Pi,w,z)$, the renormalized null integrates the
copies out, so that agent $i$ has the prior $\mathcal N(0,v_i)$ with
$v_i=s_i^2+\tau^2+\sigma^2$ in the Gaussian instance, and the frozen-environment
control updates the copy of agent $i$ from $(Q_i,Q_{-i}^{(0)})$. Along each
control the selector is $\hat\Pi$ at the control's agent laws.

\lemmaheading{The first event is a null event}{gf:first-event}
In the model \cref{gf:model}, run the label-preserving null with the
recognition family and schedule of \cref{gf:sweep}. While the structural
recognition of the full flow is $\Pi_\perp$, the full sweep and the null
sweep act identically on the agent laws. Consequently, if
$\hat\Pi(x_0)=\Pi_\perp$ and $n_1$ is the first sweep with
$\hat\Pi(x_{n_1})\neq\Pi_\perp$, the two flows have the same agent laws for
$n\leq n_1$, and $H$ and $H^0_{\mathrm{lp}}$ agree up to and including the
first event.
\begin{proof}
Under $\Pi_\perp$, \cref{gf:prior} is the product over agents of the chains
$r(dw_i)k(w_i,dz_i)p_i(y_i\mid z_i)$. In the private copy of agent $i$ every
structure places $i$ in some block with the same chain, because $r$, $k$ and
$p_i$ do not depend on the block, so the copy's block minimizer, whatever
structure it commits to, gives agent $i$ the same expected child kernel as
$\Pi_\perp$ in the full flow. Induction over the sweeps gives the claim, and
the selector is the same function of the agent laws along both flows.
\end{proof}

Hence, under block-homogeneous kernels and without multi-agent records, (F5)
against the label-preserving null can only be decided after the first event,
and a construction of this class with one selection stage fails (F5). In the
Gaussian instance the null sweep is $m_i'=a_io_i+b_ik_im_i$ with
$k_i=(\tau^2+\sigma^2)/v_i$, the renormalized null reaches
$m_i=\kappa_io_i$, $\kappa_i=v_i/(v_i+\rho_i^2)$, in one sweep, and both nulls
have the stationary means $\kappa_io_i$.

\section{The report}
\label{gf:sec-report}

\definitionheading{Declared data}{gf:declared}
Take $N=4$, $s_i^2=1$, $\rho_i^2=\tfrac14$, $\tau^2=\tfrac1{10}$, $\sigma^2=4$,
the catalog of all sixty nested structures on four agents, and the uniform
prior. Then $a_i=\tfrac45$, $b_i=\tfrac15$, $\eta_i=\tfrac15$, $k_i=\tfrac{41}{51}$
and $\kappa_i=\tfrac{102}{107}$. Write $A=\{12,34\}$, $\Pi_{\mathrm{sep}}=(A,
\{\{12\},\{34\}\})$ and $\Pi_{\mathrm{join}}=(A,\{\{12,34\}\})$. The chart is the
parameter manifold $(m_i,S_i)\in\R\times(0,\infty)$ of the Gaussian product
family, which the sweep preserves after its first step, together with the
record coordinates $o\in\R^4$; the initial structural recognition is
$\Pi_\perp$ and does not affect the dynamics, because the first sweep
replaces $Q_\nu$ by its block minimizer. The admitted class $U$ is an open
neighborhood of $(m^{(0)},o)=((-9.2,-3.0,4.2,10.4),(2.9,3.1,4.4,4.6))$, with
arbitrary initial variances, on which the conclusions of
\cref{gf:instance} hold. The selector family is the single derived map
$\hat\Pi$, so $\mathsf C$ has no real constant. For every block $I$ of every
partition at every level, $A_I$ consists of the protocol that averages one
fresh sample from each member with weights $s_i^{-2}/\lambda_I$ (for a block
of parents, from each member parent's law), and $B_I$ adds to $A_I$ and to
every $A_J$ of a block $J$ merged into $I$ one fresh sample from a named
member through the identity channel. The counterfactual class $\mathcal C$
is built from the class $\mathcal U$ of configurations with the records of
$U$ and arbitrary initial agent laws with finite second moments: its
counterfactual handoffs are the retentions \cref{gf:retention} at the actual
handoff sweep of the trajectories from $\mathcal U$ along which $\hat\Pi$ has
not reached first level $A$, restricted to the domain of the second-stage
selector, the parent states at which \cref{gf:parent-sweep} selects
$\Pi_{\mathrm{sep}}$. The records are held in $\mathcal U$ because they are
evidence and not a product of the first stage; a class that varied them
would decide (F6b) by the records.

\propositionheading{The declared instance}{gf:instance}
For the data of \cref{gf:declared} the following hold on a nonempty open set
$U$ around the declared center.
\begin{enumerate}
\item[(i)] The full history is $\Pi_\perp$ at sweep $0$, $\Pi_{\mathrm{sep}}$
at sweep $1$, and $\Pi_{\mathrm{join}}$ from sweep $2$ on, and every
relation persists by \cref{gf:persistence}; the limit is
$m^*_{\Pi_{\mathrm{join}}}=C_{\Pi_{\mathrm{join}}}(C_{\Pi_{\mathrm{join}}}+R)^{-1}o$.
After the first event the within-block contrasts equal $a_i(o_i-o_j)$.
\item[(ii)] The label-preserving null has the same partition sequence, with
the same event sweeps. The renormalized null has $\Pi_\perp$ at sweep $0$ and
$\Pi_{\mathrm{join}}$ from sweep $1$ on. Both nulls have within-block
contrasts tending to $\kappa_i(o_i-o_j)$.
\item[(iii)] The parent sweep \cref{gf:parent-sweep} with parent catalog
$\{\Pi_{\mathrm{sep}},\Pi_{\mathrm{join}}\}$ reproduces the fine block means
and selections from the handoff on; $\Pi_{\mathrm{join}}$ is its only
self-consistent structure, and from every parent state at which it selects
$\Pi_{\mathrm{sep}}$, with the records of $U$, it selects $\Pi_{\mathrm{join}}$
after finitely many sweeps.
\item[(iv)] After the second stage the retention
$(\hat m_{12}+\hat m_{34})/2$ of the union, with its record mean, satisfies
(L) to (C) with a linear parent update.
\end{enumerate}
\begin{proof}
Each history is determined by finitely many strict inequalities between the
scores \cref{gf:gaussian-score} at sweeps $0$ to $2$, and by a ball of
\cref{gf:persistence} around the limit that the trajectory enters at
sweep $3$. The scores, the trajectories, the limits, and the radius bound
obtained from the quadratic expansion of the score differences are continuous
in $(m^{(0)},o)$, so the conditions that hold at the center hold on a
neighborhood. At the center they were evaluated in double precision: the
smallest selector gap along the full and the label-preserving histories is
$0.0142$ and along the renormalized one $0.0398$, against rounding errors
below $10^{-12}$; the certificate radii in the Euclidean norm are $0.114$ for
the full flow, whose fixed-structure contraction factor is $0.188$, and
$0.261$ for the label-preserving null, whose factor is $b_ik_i=0.161$. The
first event happens at sweep $1$ for both flows, with identical states, as
\cref{gf:first-event} requires, and the null then passes through
$\Pi_{\mathrm{join}}$ on its way to $\kappa o$. For (iii), (D) of
\cref{gf:inheritance} gives the reproduction, which also holds numerically
to $10^{-15}$. The fixed point of the parent sweep under $\Pi_{\mathrm{sep}}$
is $(2.9206,4.3810)$, where $\Pi_{\mathrm{join}}$ is preferred with gap
$1.28$; under $\Pi_{\mathrm{join}}$ it is $(3.0756,4.3170)$, where
$\Pi_{\mathrm{join}}$ is preferred with gap $1.53$. From a parent state
selecting $\Pi_{\mathrm{sep}}$ the parent sweep contracts toward the first
point with factor at most $b_i$ until the selection changes, and the
selection changes at the latest when the state enters the neighborhood of
that point on which $\Pi_{\mathrm{join}}$ is preferred. For (iv), the two
parents have equal precision, so $K^c_{\Pi_{\mathrm{join}}}$ is symmetric and
has $(1,1)$ as an eigenvector, and the average of the parent means obeys an
autonomous affine contraction. The inequalities and the reproduction are
checked by the scripts recorded with the construction.
\end{proof}

The report follows the recipe of \cref{inh:recursive-formation}. The
objective is the global free energy \cref{gf:free-energy}; the dynamics is
the sweep of \cref{gf:sweep} on the committed family; the chart is that of
\cref{gf:declared}. (F1) holds in the full grade: $U$ is open in the chart and
initially unpartitioned, since $\hat\Pi=\Pi_\perp$ at sweep $0$. (F2) holds:
the selector is derived, with no constant, it follows the readout convention
of \cref{inh:null-flow} because along every control it is the original
structure minimizer at the control's agent laws, the handoffs are the first
sweeps at which the first and the second level change, which are stopping
times, and persistence is an invariant block relation. (F3) holds for the
first stage by \cref{gf:inheritance} on the set $\Omega_A$ of post-handoff
configurations from $U$, invariant by \cref{gf:instance}(i), and for the
second by \cref{gf:instance}(iv); the rules $A_I$ and $B_I$ are functions of
the index set and the constitutive data.

(F4) holds. Take $x,x'\in U$ with the same block means of $m^{(0)}$ and of
$o$ and different within-block contrasts of the records. Under $\Pi_\perp$ the
block means evolve by $\hat m_B'=a\hat o_B+bk\hat m_B$, and under structures
with first level $A$ by \cref{gf:parent-sweep}, so $R_A(x_n)=R_A(x'_n)$ for
every post-handoff sweep, while the named-member probe has the laws
$\mathcal N(m_i,\eta_i)\neq\mathcal N(m_i',\eta_i)$. The separation persists:
from the first sweep after the handoff on, the contrast is $a(o_i-o_j)$ and does
not decay. The fibers are not orbits of a symmetry of the fine field: a
translation of the within-block mean contrasts at fixed records does not
commute with the sweep, which resets them to $a(o_i-o_j)$ and makes the part
of the separation that comes from $m^{(0)}$ transient. The joint translation
$(o,m)\mapsto(o+\delta,m+a\delta)$ by a within-block contrast $\delta$
commutes with the sweep and preserves the responses of $A_I$ on the region
where the selected first level is $A$, so the persistent part is of symmetry
type there; it is not a symmetry of the selector on the whole chart, since it
changes the scores of structures with another first level. If the records
were held as constitutive constants, only the transient part would remain,
and (F4) would hold with a separation that vanishes after one sweep.

(F5) fails. Against the label-preserving null, $P=P^0_{\mathrm{lp}}$ at every
point of $U$ by \cref{gf:instance}(ii). Against the renormalized null the
sequences differ by the intermediate structure $\Pi_{\mathrm{sep}}$, which
$P^0_{\mathrm{rn}}$ skips; this is a difference under the literal comparison
of partition sequences, although it consists in the coincidence of two event
sweeps. The verdict is split in either reading, and the clause fails; since
$\mathsf C$ has no constant, there is no subset on which it holds. The
selector-free observable is the contraction of every within-block contrast
by the factor
\begin{equation}
 \frac{a}{\kappa}=\frac{s^2(v+\rho^2)}{(s^2+\rho^2)v}=\frac{214}{255}
 \label{gf:contraction}
\end{equation}
relative to both nulls, together with the different limit points
$m^*_{\Pi_{\mathrm{join}}}=(2.9956,3.1556,4.2370,4.3970)$ and
$\kappa o=(2.7645,2.9551,4.1944,4.3850)$: the interaction changes the agents'
laws but not which structures the selector reads. The frozen-environment
control was computed over $120$ sweeps at the center and gave the same
partition sequence with the same event sweeps; this is a numerical
observation at one point without a tail certificate. The strong grade does
not apply, and the amplifying grade fails against the label-preserving null
at the first event, where the two states coincide and the ratio of
\cref{inh:recursive-formation} is one.

(F6a) holds: the parents unite under the inherited sweep, by the same
selector with the same constitutive data, without resetting their records,
by \cref{gf:instance}(iii). (F6b) fails: by \cref{gf:instance}(iii) the union
forms from every counterfactual handoff in $\mathcal C$, so no subset with
nonempty interior is free of the second-stage event. The falsifier of
counterfactual dependence was tested directly. The instantaneous
second-level structure factor at the handoff depends on the parent state
$\hat m$, so it does differ between handoffs; the second-stage event and the
stationary second-level factor at the fixed point \cref{gf:fixed-point} of
the parent depend only on the parent records $\hat o$, so they are the same
for every handoff in $\mathcal C$. The second-level outcome is therefore
independent of the first-level handoff on all of $\mathcal C$, and the other
falsifier, failure of the parent to close, does not occur. With first level
$A$ the parent of \cref{gf:type-closure} is type-closed and acts on its
children by \cref{gf:top-down}.

The classification is a two-stage selection with exact inheritance: (F1),
(F2), (F3), (F4) and (F6a) hold, and (F5) and (F6b) fail. The selected
structure is also the maximum a posteriori structure of the records, but the
structure posterior is nearly flat: at the limit the structure factor
\cref{gf:structure-posterior} gives $\Pi_{\mathrm{join}}$ the mass $0.052$, and
the posterior of the records gives it $0.050$, among sixty structures. The
committed family turns a small difference of scores into a committed
selection, which the structured family would not do.

\remarkheading{Interaction-dependent selection along a regrouping path}{gf:regrouping}
Clause (F5) can be met in this model only after the first event, and only
where the full flow and the label-preserving null come to rest at different
self-consistent structures, which requires several self-consistent structures
for the same records. With $\tau^2=\tfrac1{20}$, the other data of
\cref{gf:declared}, records $(2.9,3.1,5.65,5.85)$ and initial means
$(-4,2,-9,-1)$, a computation gives the full sequence $\Pi_\perp$,
$(\{1,234\},\{1234\})$, $\Pi_{\mathrm{join}}$, the label-preserving sequence
$\Pi_\perp$, $(\{1,234\},\{1234\})$, $\Pi_{\mathrm{sep}}$, and the renormalized
sequence $\Pi_\perp$, $\Pi_{\mathrm{sep}}$. The final partitions differ from
both nulls, but the first stage selects a block that is later regrouped, so
the second stage is not a union of first-stage parents and (F6a) fails along
this path. This is a computation at one point; neither openness nor the
absence of nested paths of this kind is proved.

\openproblemheading{Formation under the global free energy}{gf:open}
Within the class $\mathfrak G$, determine whether there is an open set of
configurations on which the full flow passes through nested stages, the
first-level blocks and then their union, while the label-preserving null
comes to rest at a different structure, and on which the parent is
multistable so that the union depends on the handoff state. By
\cref{gf:first-event} such a set must lie in the regime of several
self-consistent structures, and the first event cannot contribute. Outside
$\mathfrak G$, the first event ceases to be a null event when the block
priors depend on the block or when multi-agent records are present, and the
dynamics of the structured family, in which the structure is not committed,
is open. Dynamical inheritance for child and record kernels whose update is
not affine in a record statistic needs either an augmented finite parent
state or a different notion of parent, and neither is constructed here.
